\documentclass[10pt,reqno]{amsart}

\usepackage[T1]{fontenc}
\usepackage{lmodern}
\usepackage{microtype}
\usepackage{mathtools}
\usepackage{amssymb}
\usepackage{array}
\usepackage{xcolor}
\usepackage{hyperref}

\hypersetup{
  colorlinks=true,
  linkcolor=blue!50!black,
  citecolor=blue!50!black,
  urlcolor=blue!50!black,
  pdftitle={A recursive Buchstab switching improvement for the greatest prime factor of n squared plus 1},
  pdfsubject={A greatest-prime-factor weak form of Landau's fourth problem},
  pdfkeywords={Landau's fourth problem, greatest prime factor, Harman sieve, Buchstab identity, computer-assisted proof}
}

\newcommand{\Pplus}{P^{+}}
\newcommand{\eps}{\varepsilon}
\newcommand{\Acal}{\mathcal{A}}
\newcommand{\Bcal}{\mathcal{B}}
\newcommand{\1}{\mathbf{1}}

\theoremstyle{plain}
\newtheorem{theorem}{Theorem}[section]
\newtheorem{proposition}[theorem]{Proposition}
\newtheorem{lemma}[theorem]{Lemma}
\newtheorem{corollary}[theorem]{Corollary}

\theoremstyle{remark}
\newtheorem{remark}[theorem]{Remark}

\title[Recursive Buchstab switching]{A recursive Buchstab switching improvement for the greatest prime factor of $n^2+1$}

\author{JinWen Li}
\address{SouthWest Petroleum University}
\email{lifesize1@163.com}
\date{29 August 2026}

\subjclass[2020]{11N32, 11N35, 11N75}
\keywords{greatest prime factor, quadratic polynomial, Harman sieve, Buchstab identity, Landau's fourth problem, computer-assisted proof, Lean}

\begin{document}
\raggedbottom

\begin{abstract}
We recursively expose the one-prime Buchstab tail on
$7/6<\alpha<5/4$, treating its children with the Grimmelt--Merikoski
Type~I/II estimates and the dimension-one linear sieve.  The resulting
switching saves at least $0.032303187971$ on the normalized sieve scale.
Granting these analytic inputs and the finite Mellin--Perron transfer,
infinitely many $n$ satisfy $\Pplus(n^2+1)>n^{1.323}$.  This is a
greatest-prime-factor weak form of Landau's fourth problem; the parity
barrier remains.
\end{abstract}

\maketitle

\section{Introduction}

Landau's fourth problem asks whether $n^2+1$ is prime for infinitely many
positive integers $n$.  Iwaniec proved that $n^2+1$ is a prime or a product
of two primes for infinitely many $n$ \cite{Iwaniec1978}; the possibility
that all surviving values are semiprimes is the classical parity
obstruction.  A standard quantitative complement is to make
$\Pplus(n^2+1)$ as large as possible.  Merikoski proved that
$\Pplus(n^2+1)>n^{1.279}$ infinitely often and obtained $1.312$ conditionally
on Selberg's eigenvalue conjecture \cite{Merikoski2020}.  Pascadi obtained the
unconditional exponent $1.3$ \cite{Pascadi2026}.  Grimmelt and Merikoski then
recovered $1.312$ unconditionally by proving uniform Type I and Type II
estimates for roots of quadratic congruences
\cite{GrimmeltMerikoski2025}.  Li announced $1.317$, using sharper numerical
bounds for the same Buchstab decomposition \cite{Li2025}.

These exponents quantify progress toward Landau's problem: a bound
$\Pplus(n^2+1)>n^\vartheta$ leaves a cofactor
$m<n^{2-\vartheta+o(1)}$.  Theorem~\ref{thm:main} lowers this exponent to
$0.677$, while still leaving the parity distinction $m=1$ versus
$m>1$ unresolved.

Our new input is a recursive Buchstab switching in the range where the old
argument used the linear sieve.  Merikoski already observed that Buchstab's
identity could be applied more times there to generate additional Type II
sums \cite[Remark 12]{Merikoski2020}.  With the stronger range now available
unconditionally, this produces a numerically substantial saving.  In
addition to retaining Type II three-prime children, we use the lower linear
sieve for the complementary children and recursively upper-bound the
two-prime descendants of the previously discarded one-prime tail.
Li explicitly keeps the same sieve decomposition and changes only the
numerical estimates for its terms \cite[Lemmas 2.1 and 2.2]{Li2025}; in particular,
the contribution on $7/6<\alpha<5/4$ remains the old linear-sieve term minus
the one-prime Type II term.  The three-prime term below is therefore absent
from the previously announced $1.317$ argument.

\begin{theorem}\label{thm:main}
Assume the analytic inputs in Sections~\ref{sec:transfer}
and~\ref{sec:switching}.  Specifically, assume the Type~I/II estimates of
Grimmelt and Merikoski, the dimension-one linear sieve, and the finite
Mellin--Perron transfer.  Then there are
infinitely many positive integers $n$ such that
\[
  \Pplus(n^2+1)>n^{1323/1000}=n^{1.323}.
\]
\end{theorem}

\begin{remark}[Proof-status boundary]
The hypotheses are the analytic inputs named in the theorem; they are not
formalized in Lean.  Lean checks only finite algebra, proof-adverse cell
geometry, the integer certificate, endpoint arithmetic, and the real-power
conversion.  It neither verifies those analytic inputs nor solves Landau's
fourth problem.
\end{remark}

\begin{remark}[Comparison with an isolated larger claim]
Carella's arXiv note \cite[Section 5, (5.3)--(5.4)]{Carella2025} claims the
exponent $3/2$, but the
transition from its (5.3) to (5.4) replaces
\[
 \sum_{d\mid m}\Lambda(d)=\log m
 \quad\hbox{by}\quad
 \sum_{p\mid m}\log p=\log\operatorname{rad}(m)
\]
while retaining a lower bound.  The latter quantity is at most the former,
and $n^2+1$ need not be squarefree (for example,
$7^2+1=2\cdot5^2$).  That step therefore does not establish the claimed
bound.  Our comparison frontier is the Grimmelt--Merikoski theorem and
Li's subsequent $1.317$ calculation.
\end{remark}

The gain over $1.317$ is $0.006$.  The argument is computer assisted only in
the evaluation of an explicit finite family of piecewise-rational integrals.
The certificate uses no floating-point operations.  We retain Merikoski's
original, relatively loose upper bounds for all positive Buchstab deficiencies
before $7/6$, so Li's Mathematica integrations are not an input.

\section{Transferred arithmetic information}\label{sec:transfer}

We use the notation of \cite[Section 2]{Merikoski2020}.  Let $b$ be a fixed
nonnegative smooth function supported on $[x,2x]$, put
\[
  |\Acal_d|=\sum_{\ell^2+1\equiv0\pmod d}b(\ell),
  \qquad X_b=\int_{\mathbb R}b(t)\,dt,
\]
and let
\[
  \rho(d)=\#\{\nu\pmod d:\nu^2+1\equiv0\pmod d\}.
\]
The model sequence replaces $|\Acal_d|$ by $X_b\rho(d)/d$.  For a smooth
modulus weight $\psi_P$ supported on $[P,4P]$, write
\[
  S(x,P)=\sum_p\psi_P(p)|\Acal_p|\log p.
\]

\begin{proposition}[Grimmelt--Merikoski transfer]\label{prop:transfer}
Fix a sufficiently small interior margin $\eta>0$.  The Type I estimate used
in \cite[Proposition 3]{Merikoski2020} is available with
$D\le x^{1/2-\eta}$.  If $MN=x^\alpha$ and $M\ge N$, the Type II estimate in
\cite[Proposition 4(i)]{Merikoski2020} is available, with a power-saving
error, throughout every closed interior subrange of
\begin{equation}\label{eq:typeII}
  x^{\alpha-1}<N<x^{(2-\alpha)/3},
\end{equation}
for divisor-bounded coefficients whose $N$-coefficient is supported on
squarefree integers.
\end{proposition}

\begin{proof}
Take $a=h=1$ in Grimmelt--Merikoski.  The paragraph immediately following
their Theorem~1.1 on PDF p.~2 says that, when $ah\le x^\varepsilon$, its
auxiliary prime-density hypothesis follows unconditionally from the
classical zero-free region (a possible exceptional zero only decreases the
relevant sum).  Thus no character hypothesis remains for the present fixed
shift.

Corollary~7.1 on PDF p.~22 is stated for $D\le x^{1/2-\eta}$ and gives the
sum over $d\le D$ of the absolute value of the complete smooth modulus-$d$
discrepancy.  Its test functions are normalized on $[1,2]$ and $[-1,1]$,
with constants controlled by finitely many derivative seminorms.  Every
dyadic piece used here is obtained by rescaling one fixed smooth partition,
so those seminorms, and hence the implied constant, are uniform in the
piece.  The $O((\log x)^L)$ pieces and Mellin kernels below cost only a
power of $\log x$.

Corollary~7.2 on PDF pp.~22--23 is stated for divisor-bounded coefficients,
with the $N$ coefficient squarefree-supported, and gives
\[
  x^{\alpha-1+2\eta}\le N
  \le x^{(2-\alpha)/3-4\eta/3}.
\]
In every application below, $N$ is the product of a selected subset of at
most three distinct prefix primes.  Hence it is squarefree and has at most
$\tau_3(N)$ ordered representations.  The complementary primes, the
Möbius variable, and the residual cofactor give an $M$ coefficient bounded
by $\tau_5(M)$.  The localization factors have bounded real part and unit
modulus on vertical lines, so these same fixed divisor bounds survive the
finite Mellin--Perron separation.  Thus the localized sums satisfy the
coefficient hypotheses of Corollary~7.2 term by term.  Also $N\le x^{1/3}$,
so $M=x^\alpha/N\ge N$ in all ranges used here.  Since $\eta$ is fixed and
positive, the displayed shifted interval is a closed interior subrange of
\eqref{eq:typeII}; conversely every fixed closed interior subrange is
obtained by taking $\eta$ small enough.

More explicitly, the error exponents in Corollaries 7.1 and 7.2 are,
up to an $x^{o(1)}$ factor,
\[
  1-(1-2\theta)\eta+\varepsilon/4,
  \qquad
  1-(1-2\theta)\eta+\varepsilon/8.
\]
Unconditionally $\theta=7/64$.  After fixing the sieve margin, choosing the
auxiliary $\varepsilon=\eta$ gives respectively
$1-17\eta/32$ and $1-21\eta/32$; the logarithmic losses below are therefore
absorbed by a fixed power saving.  The paragraph after Corollary~7.2 on
PDF p.~23 explicitly says that these are the Type I/II ranges used in the
proof of \cite[Theorem 2]{Merikoski2020}.  This checks all hypotheses needed
from the two corollaries; their proofs remain the cited external results.
\end{proof}

We also use Merikoski's Fundamental Proposition II with Selberg parameter
zero.  In the present notation it evaluates
\begin{equation}\label{eq:fundamentalII}
  \sum_{u\le x^{\xi-\eta}}\lambda_u
  S(\Acal(P)_u,x^{\gamma-\eta})
\end{equation}
by its model, for divisor-bounded $\lambda_u$, where
\begin{equation}\label{eq:parameters}
  \sigma=\frac{2-\alpha}{3},\qquad
  \gamma=\sigma-\alpha+1=\frac{5-4\alpha}{3},\qquad
  \xi=\frac32-\alpha.
\end{equation}
This follows from Proposition \ref{prop:transfer} by the proof of
\cite[Proposition 6]{Merikoski2020}.

\section{The recursive switching}\label{sec:switching}

Fix $7/6<\alpha<5/4$ and put
\[
  a=\alpha-1,\qquad \tau=\frac{\xi}{2}.
\]
Then $0<\gamma<\tau<a<\sigma$.  We suppress the harmless interior margins
in the displayed limiting formulae, but record one compatible choice here.
If $\nu>0$ is the parameter in Corollary~7.2 of
\cite{GrimmeltMerikoski2025}, put
\begin{equation}\label{eq:margin-parameters}
 \begin{aligned}
 a_\nu&=a+3\nu, & \sigma_\nu&=\sigma-2\nu,\\
 \gamma_\nu&=\gamma-4\nu,
 &\xi_\nu&=\xi-3\nu,\qquad \tau_\nu=\xi_\nu/2.
 \end{aligned}
\end{equation}
The available Type~II interval is
$[a+2\nu,\sigma-4\nu/3]$, so
$[a_\nu,\sigma_\nu]$ lies strictly inside it.  The proof of the
fundamental proposition works uniformly for $u\le x^{\xi_\nu}$ and sieve
level $x^{\gamma_\nu}$: its Type~I part has
\[
 x^{a+2\nu}x^{\xi_\nu}=x^{1/2-\nu},
\]
while the first Type~II prefix in its complementary part is smaller than
\[
 x^{a+2\nu+\gamma_\nu}=x^{\sigma-2\nu}
   <x^{\sigma-4\nu/3}.
\]
Thus the two published corollaries give exactly this margin version, not
merely the limiting endpoint.  We use the fixed value
$\nu=\nu_0=10^{-20}$ from Lemma~\ref{lem:explicit-margin}.  On the compact
range $\alpha<1.22$ used by the certificate,
$4\nu_0<1/25\le\gamma$ and $5\nu_0<1/25\le\sigma-a$, so
$0<\gamma_{\nu_0}<\tau_{\nu_0}<a<\sigma$.  All analytic estimates are taken
with this fixed positive margin; no limiting choice of $\nu$ is used in the
endpoint argument.

\subsection{Uniform localization and prefix sieving}

We record two standard uniformity facts because the recursive decomposition
uses them in combinations not written explicitly in the earlier papers.  The
first is the finite Mellin--Perron device used below.

\begin{lemma}[Finite cross-condition localization]\label{lem:localization}
Fix $R,C,J\ge1$.  Consider a dyadically localized bilinear sum with
$MN\asymp x^\alpha$, whose two coefficients are bounded by $\tau_R$, and
impose at most $J$ strict conditions
\[
 U(\boldsymbol m)<V(\boldsymbol n).
\]
Here $U$ and $V$ are positive integer-valued functions of the variables
assigned to the $M$- and $N$-coefficients, respectively, and
$U,V\le x^C$.  A fixed smooth function of $MN/P$ and a factor $\log(MN)$
may also be present.  Then the sum is a multiple Mellin--Perron integral of
ordinary Type~II sums whose coefficients are bounded by
$O_{C,J}(\tau_R)$.  The total $L^1$ mass of the kernels is
$O((\log x)^{O_{C,J}(1)})$, and the truncation error is $O_A(x^{-A})$
for any fixed $A>0$, after increasing the truncation height.
Consequently any fixed power-saving Type~II estimate that is uniform for
divisor-bounded coefficients survives all these localizations.
\end{lemma}

\begin{proof}
Since $U,V$ are integers, the strict comparison is equivalent to
$U<V-1/2$.  Put $y=(V-1/2)/U$ and $c=1/\log x$.  For $y\ne1$, the
truncated Perron formula is
\[
 \frac{1}{2\pi i}\int_{c-iT}^{c+iT}y^s\frac{ds}{s}
 =\1_{y>1}+O\!\left(
 y^c\min\left\{1,\frac1{T|\log y|}\right\}\right).
\]
Here is the required estimate, rather than an appeal to an unquantified
localization principle.  The full vertical integral equals $1$ for $y>1$
and $0$ for $0<y<1$: close the contour to the left or right, respectively,
and take the residue at $s=0$.  On either omitted vertical tail, integration
by parts in
\(e^{it\log y}/(c+it)\) gives $O(y^c/(T|\log y|))$ when
$T|\log y|\ge1$; the trivial bound for the complete transition gives
$O(y^c)$ otherwise.  This proves the displayed formula.

The half-integer cut makes the estimate uniform.  If $U<V$, then
$y\ge1+1/(2U)$ and
\[
 \log y\ge\frac{2}{4U+1}.
\]
If $U\ge V$, then $1-y\ge1/(2U)$ and
$-\log y\ge1-y\ge1/(2U)$.  Hence in both cases
\[
 |\log y|\ge\frac{2}{4U+1}\ge\frac{2}{5x^C}.
\]
Moreover $y^c\le e^C$ when $y>1$, and $y^c\le1$ when $y<1$.
Thus one comparison has pointwise truncation error
$O_C(x^C/T)$.

The replacement of $V$ by $V-1/2$ only excludes the equality case.  It
does not change the underlying representations or their divisor bounds.

For a cross-condition $U(\boldsymbol m)<V(\boldsymbol n)$, attach
$U^{-s}$ to the $M$-coefficient and $(V-1/2)^s$ to the $N$-coefficient.
If the smaller variable belongs to the $N$ side, attach $U^{-s}$ there and
$(V-1/2)^s$ to $M$; the designation of the squarefree Type-II family is not
changed.  A product cutoff has the same form:
if $A_MA_N\le Q$, with $A_M$ and $A_N$ formed from variables on the two
sides, write it as
\[
 A_M<\left\lfloor Q/A_N\right\rfloor+1.
\]
Thus $U=A_M$ is assigned to the $M$ side and
$V=\lfloor Q/A_N\rfloor+1$ to the $N$ side; the complementary inequality is
handled by subtracting this class.  This also covers any exact priority
cutoff not already fixed by the exponent cell.
If $U$ or $V$ depends on a particular prime-factor representation rather
than only on the collected integer, attach the phase before collecting that
coefficient.  The number of prime slots is fixed and the vertical factors
are uniformly bounded, so summing those representations preserves a fixed
divisor bound.

The recursive branches specialize this rule as follows.
\begin{description}
\item[\(A_3\), Type II.]
For the first eligible subset \(I\), put
\(N=\prod_{i\in I}q_i\) and
\(M=(\prod_{i\notin I}q_i)ds\).
An ordering comparison split across the families assigns the smaller prime
to \(U^{-s}\) and the larger to \((V-1/2)^s\).
If \(q_3\in N\), roughness is
\(P^+(d)<q_3\), with \(P^+(d)^{-s}\) on \(M\) and
\((q_3-1/2)^s\) on \(N\); if \(q_3\notin N\), it is internal to \(M\).
Priority product cutoffs are fixed away from cell boundaries, or are
separated by the preceding \(A_MA_N\le Q\) rule.

\item[Direct pair Type II.]
Take \(N=q_1q_2\) and put the sifted cofactor and Mobius variable on \(M\).
The order \(q_2<q_1\) is internal to \(N\); for \(d>1\), the only cross
condition is \(P^+(d)<q_2\), assigned as above.  The term \(d=1\) is automatic.

\item[Type-I prefix branches.]
The \(A_3\) lower-sieve branch has prefix \(q_1q_2q_3\) and cutoff \(q_3\);
the universal \(U_{\rm LS}\) branch has prefix \(q_1q_2\) and cutoff \(q_2\).
Their ordering is internal and the variable cutoff is carried by the
prefix-dependent Rosser weight of Proposition~\ref{prop:prefix-sieve};
neither branch invokes this Type-II Perron lemma.

\item[Fixed-cutoff bases.]
The one-prime tail base has prefix \(q_1\), and the recursive base has
prefix \(q_1q_2\), both with cutoff \(x^{\gamma_\nu}\).
Their cell conditions are fixed before the Fundamental Proposition is
applied.  The recursive children reuse the \(A_3\) assignments above; the
choice between direct and recursive upper bounds is a boxwise decision, not
a Perron condition.
\end{description}

For each priority class the number $J$ of genuine cross-comparisons is fixed,
independent of $x$.  On $\Re s=c$, the real factors are at most $e^C$ and
the imaginary factors have modulus one.  Also
\[
 \int_{-T}^{T}\frac{dt}{|c+it|}
 \ll\log(T\log x)\ll_{A'}\log x
 \qquad(T=x^{A'}).
\]
Replace the $J$ indicators one at a time and telescope.  If the unreduced
sum has at most $x^B$ tuples, the total truncation error is
\[
 O_{C,J}\!\left(
 x^B\frac{x^C}{T}(\log x)^{J-1}\right).
\]
Given $A>0$, choosing $A'>A+B+C+J+2$ makes this
$O_A(x^{-A})$.  The resulting multiple integral has total kernel mass
$O((\log x)^J)$.  Therefore a Type~II error $O(x^{1-\delta})$, uniform for
divisor-bounded coefficients, becomes
$O(x^{1-\delta}(\log x)^J)=O(x^{1-\delta/2})$ for large $x$.

For a smooth compactly supported $W$, Mellin inversion separates
$W(MN/P)$; rapid decay of its Mellin transform gives the same arbitrary
power truncation error.  Finally
$\log(MN)=\log M+\log N$, with the logarithms handled by partial
summation.  All vertical factors have size $O_C(1)$, so the coefficients
remain uniformly divisor bounded.  The fixed polylogarithmic kernel mass
is absorbed after decreasing the power-saving exponent.  This proves the
lemma.  Merikoski's use of Perron in the proof of Fundamental Proposition~II
and his Mellin removal of the smooth product weights provide the qualitative
templates \cite[Sections 2.3 and 3.4]{Merikoski2020}.  This proof is an
ordinary analytic argument and is not part of the Lean formalization.
\end{proof}

\begin{proposition}[Prefix-uniform linear sieve]\label{prop:prefix-sieve}
Let a dyadically localized ordered prefix
$\boldsymbol q=(q_1,\ldots,q_j)$ consist of at most three primes.  Put
$u=q_1\cdots q_j$ and suppose that the localized cell gives the uniform
upper bound $u\le x^{s_0}$, with $1/2-s_0-\nu>0$.  All prefix primes are at least the sieve cutoff
$z=z(\boldsymbol q)$, with $z(\boldsymbol q)\asymp x^\zeta$; in the variable
cutoff branches, $z(\boldsymbol q)=q_j$ is the least prefix prime.  The
Type~I estimate in Proposition~\ref{prop:transfer}, summed over any fixed
finite collection of such cells with divisor-bounded prefix weights, permits
the standard dimension-one upper and lower linear sieve with remainder level
\[
 D=x^{1/2-s_0-\nu}.
\]
After removal of the common Euler factor, the corresponding sieve parameter
is $s=(1/2-s_0-\nu)/\zeta$ and the main-term multipliers are
$e^{-\gamma_E}F(s)$ and $e^{-\gamma_E}f(s)$.  The lower multiplier is
understood to be zero for $s\le2$.  The aggregate remainder is
$O(x^{1-\delta})$ for some $\delta=\delta(\nu)>0$.
\end{proposition}

\begin{proof}
For each localized prefix, insert the upper or lower Rosser weights at
level \(D=x^{1/2-s_0-\nu}\).  The per-prefix inequalities are the
dimension-one formulae (12.12)--(12.14) of
\cite[Chapter 12, pp. 235--236]{FriedlanderIwaniec2010}.  Their remainders
are kept with their signs.  After multiplying by the nonnegative prefix
partition and summing, put \(k=um\) and \(e=ud\).  This gives, before any
absolute value is taken,
\[
 \sum_{\boldsymbol q}c(\boldsymbol q)R_{\boldsymbol q}^{\pm}
 =\sum_e\Lambda_e^\pm r_P(e).
\]
Here \(r_P(e)\) is exactly the smooth modulus-\(e\) discrepancy occurring
inside the absolute value in Grimmelt--Merikoski Corollary~7.1; in
particular it is independent of the prefix representation of \(e\).
Appendix~\ref{app:prefix-sieve} proves
\[
 |\Lambda_e^\pm|\ll(\log x)^B\tau_K(e),\qquad
 ud\le x^{1/2-\nu},
\]
so the absolute discrepancy sum in
\cite[Corollary 7.1, p. 22]{GrimmeltMerikoski2025}, with \(a=h=1\),
leaves a fixed power saving after all prefix cells are summed.

The appendix also evaluates the model sums uniformly.  The primes in a
Rosser modulus are below the least prefix prime, so they are coprime to
\(u\); Merikoski's Lemma~8 and the dimension-one fundamental lemma therefore
give the same Euler product as in the unprefixed sequence.  For the \(A_3\)
lower branch the three-prime level remains a positive power of \(x\) at
\(\alpha=7/6\), while \(s\le2\) contributes zero.  For the two-prime upper
branch the resulting coefficient is \(U_{\rm LS}\); if \(0<s<1\), positivity
allows the cutoff to be lowered to \(z_*=D^{1-\omega}\) before applying the
upper sieve.  Finally Merikoski's Lemma~7, with the sieve multiplier frozen
on each exponent cell, converts the model sums to \eqref{eq:A3lower} and
\eqref{eq:Tlower}.  The hypotheses retained in
Remark~\ref{rem:prefix-sieve-limit} are precisely the cited external
analytic theorems, rather than an additional cancellation assumption.
\end{proof}

\begin{remark}[Limitation of the prefix-sieve transfer]
\label{rem:prefix-sieve-limit}
Proposition~\ref{prop:prefix-sieve} is not a Lean theorem.  Its external
inputs are precisely Grimmelt--Merikoski Corollary~7.1 and the
dimension-one fundamental lemma of \emph{Opera de Cribro}; neither source
theorem is reproved here.  Appendix~\ref{app:prefix-sieve} checks their
hypotheses for the least-prefix cutoff, including the uniform Euler product,
the normalized smooth family, and the collected absolute remainder.  Thus
no additional ``prefix-uniformity'' assertion is left as an assumption.
\end{remark}

The decomposition preceding the old linear-sieve term is
\begin{equation}\label{eq:retain-F6}
 S(\Acal(P),2\sqrt P)
 \le S(\Acal(P),x^a)
 -\sum_{x^a<q<x^\sigma}S(\Acal(P)_q,q).
\end{equation}
The negative sum in \eqref{eq:retain-F6} is retained unchanged; its model
contribution is the term $F_6$ bounded in Section~5.  We improve only the
first term.  The first Buchstab split is
\begin{equation}\label{eq:first-switch}
 \begin{aligned}
  S(\Acal(P),x^a)
  &=S(\Acal(P),x^\gamma)
   -\sum_{x^\gamma<q_1<x^\tau}S(\Acal(P)_{q_1},q_1)-R_1,\\
  R_1&:=\sum_{x^\tau<q_1<x^a}S(\Acal(P)_{q_1},q_1).
 \end{aligned}
\end{equation}
Applying Buchstab twice more to the retained sum yields
\begin{equation}\label{eq:four-terms}
  S(\Acal(P),x^a)=A_0-A_1+A_2-A_3-R_1,
\end{equation}
where
\begin{align*}
 A_0&=S(\Acal(P),x^\gamma),\\
 A_1&=\sum_{x^\gamma<q_1<x^\tau}
       S(\Acal(P)_{q_1},x^\gamma),\\
 A_2&=\sum_{x^\gamma<q_2<q_1<x^\tau}
       S(\Acal(P)_{q_1q_2},x^\gamma),\\
 A_3&=\sum_{x^\gamma<q_3<q_2<q_1<x^\tau}
       S(\Acal(P)_{q_1q_2q_3},q_3).
\end{align*}
The first three terms are covered by \eqref{eq:fundamentalII}, since
$q_1\le x^\tau$ and $q_1q_2\le x^{2\tau}=x^\xi$.

For $A_3$, retain only tuples for which at least one of
\begin{equation}\label{eq:subset-sums}
 \beta_1+\beta_2,\quad \beta_1+\beta_3,\quad
 \beta_2+\beta_3,\quad \beta_1+\beta_2+\beta_3
\end{equation}
lies in the shifted Type II interval $[a_\nu,\sigma_\nu]$.  The
corresponding product of
distinct primes is squarefree and may be chosen as the Type II variable.
Use the canonical $1200\times300\times160$ exponent mesh of
Proposition~\ref{prop:certificate}.  A cell meeting a Type~II endpoint or a
priority boundary is assigned the proof-adverse lower contribution (zero if
no other lower sieve is available).  Thus its loss is already included in
the certified saving, rather than being postponed to a limiting mesh
argument.  On every remaining cell,
Lemma~\ref{lem:localization} removes the finitely many ordering and
roughness cross-conditions and leaves divisor-bounded coefficients.  This
Type~II part gives one lower bound for $A_3$.  Unlike the earlier
three-step argument, we also
lower-bound complementary triples by the dimension-one lower linear sieve
and take the larger available weight.  Since $A_3$ occurs with a negative
sign, every such lower bound may be subtracted in \eqref{eq:four-terms}.

\begin{lemma}\label{lem:A3-transfer}
After restoring fixed interior margins, the retained part of $A_3$ has the
same asymptotic formula for $\Acal(P)$ and $\Bcal(P)$, with a power-saving
error on each smooth $P$-block.
\end{lemma}

\begin{proof}
Order the four subsets in \eqref{eq:subset-sums} and assign each good prime
triple to the first subset whose exponent sum lies in the fixed interior of
$[a_\nu,\sigma_\nu]$.  For every earlier subset, nonmembership is the disjoint union
of the alternatives ``below the lower endpoint'' and ``above the upper
endpoint''.  Expanding these alternatives produces only finitely many
priority classes (at most $1+2+2^2+2^3$ before empty classes are removed).
Smoothly localize the three prime variables.  For an assigned
subset $I$, put $N=\prod_{i\in I}q_i$ and put the complementary prime factors
and the sifted cofactor into $M$.  Then $MN\asymp x^\alpha$,
\[
 x^{a_\nu}\le N\le x^{\sigma_\nu},\qquad M\ge N,
\]
and $N$ is squarefree.  The coefficient of $N$ has only divisor-boundedly
many prime representations.  The coefficient of $M$ is also divisor
bounded, since it records at most the complementary prime factors and one
cofactor.

Here are the coefficient and roughness details.  From the definition of a
sifted sum,
\[
 S(\Acal(P)_{q_1q_2q_3},q_3)
 =\sum_{(r,P(q_3))=1}|\Acal_{q_1q_2q_3r}|
   \psi_P(q_1q_2q_3r)\log(q_1q_2q_3r).
\]
Insert the finite identity
\[
 \1_{(r,P(q_3))=1}
   =\sum_{d\mid r,\,d\mid P(q_3)}\mu(d)
\]
and write $r=ds$.  If $N=\prod_{i\in I}q_i$ is the selected subset,
put $M=(\prod_{i\notin I}q_i)ds$.  After dyadic subdivision, the resulting
sum has exactly the bilinear form of Proposition \ref{prop:transfer}.
The $N$-coefficient is supported on squarefree integers: the selected
$q_i$ are distinct.  More quantitatively, if $t=|I|$, then before the
separation of cross-conditions the numbers of possible factorizations obey
the uniform bounds
\[
 \#\{N=\textstyle\prod_{i\in I}q_i\}\le \tau_t(N),\qquad
 \sum_{M=(\prod_{i\notin I}q_i)ds}|\mu(d)|
       \le \tau_{5-t}(M).
\]
The harmless conventions $\tau_0(1)=1$ and
$\tau_j\le\tau_5$ may be used.  Thus both coefficients are bounded by a
fixed divisor function.  Repeated prime factors shared by $M$ and $N$ do
not affect this conclusion: Corollary 7.2 requires squarefree support only
for the $N$ coefficient and imposes no coprimality condition on $M,N$.

There is one genuine cross-condition when $q_3$ belongs to the selected
product $N$.  For $d>1$ the condition $d\mid P(q_3)$ is equivalently
\[
 \mu(d)\ne0,\qquad P^+(d)<q_3;
\]
the term $d=1$ is automatic.  Hence the only datum from the $M$ side which
must be compared with the $N$ side is the single integer $P^+(d)$.  A
truncated Perron integral separates
$P^+(d)<q_3$ into $(P^+(d))^{it}$ and $q_3^{-it}$.  If $q_3\notin N$, this
condition is already internal to the $M$ coefficient.  This makes explicit
that expanding the roughness condition does not create an unbounded number
of cross variables.

It remains to separate conditions that involve variables on both sides.
Prime ordering and the strict comparison $P^+(d)<q_3$ have exactly the
integer form covered by Lemma~\ref{lem:localization}; using $q_3-1/2$ in
the latter comparison also removes any equality ambiguity in truncated
Perron inversion.  The selected-subset priority is already fixed on the
interior exponent cell.  The product cutoff and logarithm are handled by
the Mellin and partial-summation separations in that lemma.  Hence the two
coefficient bounds above remain $O(\tau_5(M))$ and
$O(\tau_5(N))$, while all kernel and dyadic losses are powers of
$\log x$.  These are absorbed by the fixed power saving in Proposition
\ref{prop:transfer}.

This is precisely the mechanism in the proof of Merikoski's Fundamental
Proposition and in the evaluation of its three-prime term
\cite[Sections 2.3, 2.4.3, and 3.4]{Merikoski2020}.  Most importantly, the
paper explicitly states there that \emph{any combination} of
$q_1,q_2,q_3$ lying in the Type II range may be used.  The new term changes
only the outer exponent boxes; every retained box lies in the unconditional
interior supplied by Proposition \ref{prop:transfer}.  Summing the four
disjoint priority classes proves the lemma.
\end{proof}

All conversions of model almost-prime sums below use exactly
\cite[Lemma 7]{Merikoski2020}: its rough-number input is equation~(2.4), its
measure is $d\beta_1\cdots d\beta_k/(\beta_1\cdots\beta_{k-1}\beta_k^2)$,
and the numerical Buchstab bounds used here are precisely equation~(2.5).
No stronger model-integral statement is assumed.

Let $\omega$ be the Buchstab function and set
\[
  u_+=0.5617,\qquad u_-=0.5607.
\]
The bounds
\begin{equation}\label{eq:omega-bounds}
  \omega(v)\le u_+\quad(v\ge4),\qquad
  \omega(v)\ge u_-\quad(v\ge2.47)
\end{equation}
are quoted in \cite[(2.5)]{Merikoski2020}.  On the retained range
$\alpha<1.22$, an ordered three-prime child has
\[
 \frac{\alpha-\beta_1-\beta_2-\beta_3}{\beta_3}
 >\frac{3-2\alpha}{\alpha-1}>2.54.
\]
The one- and two-prime fundamental-proposition bases have arguments at
least $9$ and $7.5$, respectively.  If $\beta_1+\beta_2\le\sigma$, the
ordered Type II pair has argument at least
$2(\alpha-\sigma)/\sigma\ge6.4$.  Thus every use of
\eqref{eq:omega-bounds} is within its stated range.

We also use the standard dimension-one linear-sieve functions $F$ and $f$
\cite[Chapter 12]{FriedlanderIwaniec2010}.  Define the rational lower
envelope
\begin{equation}\label{eq:Phi}
 \Phi(s)=
 \begin{cases}
  0,&s<2,\\
  4(s-2)/s^2,&2\le s\le4,\\
  1/2,&s\ge4.
 \end{cases}
\end{equation}
For $2\le s\le4$ the identity
$e^{-\gamma_E}f(s)=2\log(s-1)/s$ and
$\log(1+t)\ge2t/(2+t)$ give
\begin{equation}\label{eq:Phi-lower}
 \Phi(s)\le e^{-\gamma_E}f(s).
\end{equation}
At $s=4$ the rational value is $1/2$, and the standard monotonicity of $f$
gives the continuation in \eqref{eq:Phi}.  Similarly,
$e^{-\gamma_E}F(s)=2/s$ for $1\le s\le3$ and is at most $2/3$ for
$s\ge3$.

For an ordered triple put
\[
 s_3=\frac{1/2-\beta_1-\beta_2-\beta_3}{\beta_3},
 \qquad
 W_3=\max\{u_-\1_{\rm good},\Phi(s_3)\},
\]
where $\1_{\rm good}$ indicates that one of \eqref{eq:subset-sums} lies in
$[a,\sigma]$.  Proposition \ref{prop:transfer} handles the first weight;
Proposition~\ref{prop:prefix-sieve} handles the second.  Indeed, after a
fixed prefix of total exponent $s_0$, the available sieve level is
$x^{1/2-s_0-\nu}$, and combining its Rosser weight with at most three prime
indicators still gives a fixed divisor-bounded coefficient.  Hence
\begin{equation}\label{eq:A3lower}
 A_3^{\rm low}(\alpha)=\alpha
 \int_{\gamma<\beta_3<\beta_2<\beta_1<\tau}
 W_3\,\frac{d\beta_1\,d\beta_2\,d\beta_3}
 {\beta_1\beta_2\beta_3^2}
\end{equation}
is a valid limiting lower model for $A_3$.

It remains to retain part of $R_1$.  For $\tau<\beta_1<a$, Buchstab gives
\[
 S(\Acal(P)_{q_1},q_1)
 =S(\Acal(P)_{q_1},x^\gamma)
  -\sum_{x^\gamma<q_2<q_1}S(\Acal(P)_{q_1q_2},q_2).
\]
Put $\delta_2=1/2-\beta_1-\beta_2$ and define the proof-adverse upper
linear-sieve envelope
\[
 U_{\rm LS}(\beta_1,\beta_2)
 =\max\left\{\frac2{\delta_2},\frac2{3\beta_2}\right\}.
\]
For $\beta_1+\beta_2\le\xi$, put
\[
 C_{\rm rec}(\beta_1,\beta_2)
 :=\frac{u_+}{\gamma}
   -\int_\gamma^{\beta_2}
      W_3(\beta_1,\beta_2,\beta_3)\frac{d\beta_3}{\beta_3^2}.
\]

\begin{lemma}[Nonnegativity of the recursive candidate]
\label{lem:recursive-candidate-nonnegative}
On its stated domain, $C_{\rm rec}(\beta_1,\beta_2)\ge0$.
\end{lemma}

\begin{proof}
Apply Buchstab's identity to the model sum with prefix $q_1q_2$ and cutoff
$x^\gamma$.  Its base is the sum of the children with
$x^\gamma<q_3<q_2$ and the remaining nonnegative $q_2$-rough term.  Denote
their normalized model contributions by $B$, $C$, and $R$, respectively;
then $B=C+R$ with $R\ge0$.  Fundamental Proposition~II and
\eqref{eq:omega-bounds} give the base upper bound
$B\le u_+/\gamma$, while \eqref{eq:A3lower} gives
\[
 C\ge\int_\gamma^{\beta_2}
 W_3(\beta_1,\beta_2,\beta_3)\frac{d\beta_3}{\beta_3^2}.
\]
Therefore
\[
 C_{\rm rec}\ge B-C=R\ge0.
\]
This is an inequality between the analytic model terms; the finite program's
explicit nonnegativity guard is only a proof-adverse rounding safeguard.
\end{proof}

Let $U_2(\beta_1,\beta_2)$ be the minimum of $U_{\rm LS}$ and the following
available upper bounds:
\begin{align*}
 &\frac{u_+}{\beta_2},
 &&\text{if }\beta_1+\beta_2\in[a,\sigma],\\
 &C_{\rm rec}(\beta_1,\beta_2),
 &&\text{if }\beta_1+\beta_2\le\xi,
\end{align*}
The second line is one more Buchstab identity: Fundamental Proposition II
evaluates its base and \eqref{eq:A3lower} lower-bounds its children.
The proof-adverse linear-sieve option is uniform over the two-prime prefix
by Proposition~\ref{prop:prefix-sieve}.
Consequently
\begin{equation}\label{eq:Tlower}
 T(\alpha)=\alpha\int_\tau^a
 \left[\frac{u_-}{\gamma}
 -\int_\gamma^{\beta_1}U_2(\beta_1,\beta_2)
       \frac{d\beta_2}{\beta_2}\right]_+
 \frac{d\beta_1}{\beta_1}
\end{equation}
is a lower model for $R_1$.

\subsection{Complete branch-to-input ledger}

Put
\[
 \mathfrak M(P):=X_b\int_0^\infty\psi_P(u)\,\frac{du}{u}.
\]
This is the common smooth block factor in
\cite[Lemma 7]{Merikoski2020}.  For a fixed $\nu>0$ and a fixed exponent
mesh, let $H_\nu(\alpha)$ denote the expression below with the shifted
parameters in \eqref{eq:margin-parameters} and the proof-adverse cell
endpoints used by the certificate.

\begin{proposition}[Analytic branch ledger]\label{prop:branch-ledger}
Uniformly for $7/6<\alpha<1.22$, the recursive decomposition uses exactly
the one-sided estimates in the following table.  Each row is valid with a
power-saving remainder before the finite sums of dyadic and Perron pieces.

\begin{center}
\small
\renewcommand{\arraystretch}{1.12}
\begin{tabular}{|>{\raggedright\arraybackslash}p{0.13\textwidth}|
>{\raggedright\arraybackslash}p{0.29\textwidth}|
>{\raggedright\arraybackslash}p{0.43\textwidth}|}
\hline
term & analytic input & proof direction and condition \\
\hline
$A_0$ & Fundamental Proposition II & upper bound, $u=1$, using
$u_+/\gamma$ \\
\hline
$A_1$ & Fundamental Proposition II & lower bound on the common one-prime
domain, using $u_-/\gamma$ \\
\hline
$A_2$ & Fundamental Proposition II & upper bound on a containing ordered
two-prime domain, using $u_+/(2\gamma)$ after symmetry \\
\hline
$A_3$ & Corollary 7.2 & lower bound whenever a complete subset-product cell
lies in $[a_\nu,\sigma_\nu]$ \\
\hline
$A_3$ & Proposition~\ref{prop:prefix-sieve} & lower linear-sieve bound on
the same child; take the larger of the two available lower bounds \\
\hline
$R_1$ base & Fundamental Proposition II & lower bound because
$\beta_1<a<\xi$ \\
\hline
two-prime child & Proposition~\ref{prop:prefix-sieve} & universal upper
linear-sieve envelope $U_{\rm LS}$ \\
\hline
two-prime child & Corollary 7.2 or one more Buchstab split & take the smaller
upper bound: direct Type~II when $\beta_1+\beta_2\in[a_\nu,\sigma_\nu]$;
otherwise base upper minus $A_3$ lower when
$\beta_1+\beta_2\le\xi_\nu$ \\
\hline
\end{tabular}
\end{center}

Consequently, in the normalization above,
\[
 S(\Acal(P),x^a)
 \le \bigl(H_\nu(\alpha)+o(1)\bigr)\mathfrak M(P),
\]
where the $o(1)$ is uniform on the fixed shifted cells.  The quantitative
comparison between this fixed-margin envelope and $H(\alpha)$ is given
after Proposition~\ref{prop:certificate}; no unquantified mesh limit is used
in the endpoint budget.
\end{proposition}

\begin{proof}
The signs $A_0-A_1+A_2-A_3-R_1$ determine the required directions.  The
first three rows follow from \eqref{eq:fundamentalII}, since their prefix
exponents are at most $0,\tau$, and $2\tau=\xi$, respectively.  The two
$A_3$ rows are Lemma~\ref{lem:A3-transfer} and
Proposition~\ref{prop:prefix-sieve}; the maximum of two lower bounds for the
same nonnegative child remains a lower bound.

For $R_1$, Buchstab's identity writes each node as a base minus its
two-prime children.  A base lower bound minus children upper bounds is
therefore a node lower bound.  Each child has the universal upper
linear-sieve bound.  On a complete Type~II pair cell, Corollary 7.2 supplies
the second upper bound.  If the pair exponent is at most $\xi_\nu$, one
more Buchstab identity instead gives a base upper bound minus the already
proved three-prime lower bound.  Taking the minimum of valid upper bounds
is valid.  Lemma~\ref{lem:localization} preserves all these estimates after
the finitely many cross-condition separations.  Summing the rows proves the
claim.
\end{proof}

\begin{remark}[Exhaustive branch coverage]\label{rem:branch-coverage}
The supplementary branch-coverage audit checks the complete canonical grid
using integer comparisons only.  It visits $34\,215\,168$
ordered $A_3$ boxes and $19\,635\,200$ two-prime tail boxes.  Every row of
the ledger occurs.  The proof-adverse minimum Buchstab arguments are,
respectively,
\[
 5.206977117,\qquad 9.007506255,\qquad
 8.002880298,\qquad 7.508771632
\]
for the $A_3$ Type~II lower branch, the $R_1$ base, the direct pair
Type~II upper branch, and the recursive-base upper branch.  Thus the first
two exceed $2.47$ and the last two exceed $4$, as required for the constants
$u_-$ and $u_+$.  The audit terminates with
\texttt{BRANCH\_LEDGER\_CERTIFIED=YES}.
\end{remark}

Finally put $L=\log(\tau/\gamma)$.  The complete recursive upper envelope is
\begin{equation}\label{eq:H}
 H(\alpha)=\frac{\alpha}{\gamma}
 \left(u_+-u_-L+\frac{u_+}{2}L^2\right)
 -A_3^{\rm low}(\alpha)-T(\alpha).
\end{equation}
The original linear-sieve envelope is $4\alpha$, so the new saving is
\begin{equation}\label{eq:Delta}
 \Delta=\int_{7/6}^{5/4}
 \max\{0,4\alpha-H(\alpha)\}\,d\alpha.
\end{equation}

\section{Integer certificate for the switching integral}

\begin{proposition}\label{prop:certificate}
The integral in \eqref{eq:Delta} satisfies
\[
  \Delta\ge\frac{32303187971}{10^{12}}
  =0.032303187971>\frac4{125}.
\]
\end{proposition}

\begin{proof}
We describe the exact certificate implemented in the accompanying file
\texttt{scripts/certify\_harman\_recursive\_tail.cpp}.  It uses $N=1200$
equal $\alpha$-cells, $B=300$ cells for \eqref{eq:A3lower}, and $C=160$
cells for the tail \eqref{eq:Tlower}.  All quantities are stored as integer
multiples of $10^{-12}$.  Products use signed $128$-bit intermediates;
lower bounds are rounded down, upper bounds are rounded up, and every
narrowing conversion is checked for overflow.

For an $\alpha$-cell $[\alpha_i,\alpha_{i+1}]$, monotonicity gives the common
domain
\[
 [\gamma(\alpha_i),\tau(\alpha_{i+1})]
\]
and the containing domain
\[
 [\gamma(\alpha_{i+1}),\tau(\alpha_i)].
\]
The first three terms in \eqref{eq:H} are bounded, in proof-adverse
directions, by
\begin{align*}
 A_{0,i}^{+}&=u_+\frac{\alpha_{i+1}}{\gamma(\alpha_{i+1})},\\
 A_{1,i}^{-}&=u_-\frac{\alpha_i}{\gamma(\alpha_i)}R^-_B,
 \\
 A_{2,i}^{+}&=\frac{u_+}{2}
 \frac{\alpha_{i+1}}{\gamma(\alpha_{i+1})}(R^+_B)^2,
\end{align*}
where $R^-_B$ is the right-endpoint lower Riemann sum for
$\int d\beta/\beta$ over the common domain, and $R^+_B$ is the left-endpoint
upper sum over the containing domain.

For the three-prime lower term, each ordered $\beta_1,\beta_2$ box produces
a finite family of rational $\beta_3$ intervals.  Some are uniformly Type
II throughout the entire $\alpha,\beta_1,\beta_2$ cell.  The others come
from the 21 rational lower-sieve thresholds between $2.01$ and $4$.
Overlapping intervals take the largest certified lower weight.  The
$\beta_3$ integral is evaluated exactly by
\[
 \int_r^s\frac{d\beta_3}{\beta_3^2}=\frac1r-\frac1s,
\]
and $1/(\beta_1\beta_2)$ is replaced by its value at the upper box
endpoints.  Thus the resulting $A_{3,i}^{-}$ is a lower bound.

For the recursive tail, a $\beta_1$ cell lies in the common intersection
$[\tau(\alpha_i),a(\alpha_i)]$.  Its child $\beta_2$ cells cover a containing
domain from $\gamma(\alpha_{i+1})$ to the upper endpoint of the
$\beta_1$ cell.  Each child upper bound is the minimum of the proof-adverse
linear-sieve bound, a uniformly available Type II bound, and the recursive
base-minus-children bound when the two-prime prefix is at most $\xi$.
The same-exponent-cell region is subtracted in full.  This deliberately
overcounts both distinct primes in one cell and the harmless region
$\beta_2\ge\beta_1$; it cannot make the tail lower bound too large.

It follows pointwise on the cell that
\[
 \max\{0,4\alpha-H(\alpha)\}
 \ge \max\{0,4\alpha_i-
 (A_{0,i}^{+}-A_{1,i}^{-}+A_{2,i}^{+}
   -A_{3,i}^{-}-T_i^{-})\}.
\]
Only $\alpha<1.22$ is retained; zero is used elsewhere.

The exact integer output is
\begin{verbatim}
scale=1000000000000
alpha_steps=1200 beta_steps=300 tail_beta_steps=160
positive_alpha_cells=727
tail_lower_integral=0.018330264130
a3_lower_integral=0.044459636636
saving_lower_units=32303187971
saving_lower=0.032303187971
target=0.032000000000
CERTIFIED=YES
\end{verbatim}
The last comparison is the integer inequality
$32303187971>32000000000$.

The module \texttt{HarmanRecursiveCertificate.lean} is a third Lean
implementation of the complete fixed-point algorithm.  Lean evaluates the
canonical grid by \texttt{native\_decide} in four 192-cell
blocks and one zero block.  The respective unscaled saving sums are
\[
 154113574025565,\quad151615832402575,\quad
 117510700363272,\quad41925800002266,\quad0.
\]
Their sum is $465165906793678$; downward division by $12N=14400$ is
$32303187971$.  Small projection modules connect every expensive block
computation to its exact natural-number value.  Thus the large block
equalities use Lean's native evaluator; after those equalities are available,
\texttt{HarmanRecursiveCertificateCanonical.lean} checks the final natural-number
sum, division, rational normalization, and strict endpoint inequalities by
kernel elaboration without re-expanding the blocks.

As an implementation-level cross-check, the accompanying Python audit
program independently
reconstructs every cell using Python arbitrary-precision integers.  It does
not call or parse the C++ executable and uses an independent sweep-line
algorithm for weighted interval maxima.  On the canonical grid it reports
the same totals.  With the optional \texttt{--dump-cells} flag, all 1200
canonical cell records agree byte for byte; both dumps have SHA-256
\par\smallskip
\noindent\hspace*{1em}\texttt{6e1901ca1dc8b2c32ae6bcc3759a42ea}\par
\noindent\hspace*{1em}\texttt{9a3dab8932870eeb3f6e2d42f0ac3175}.\par\smallskip
Thus the Lean, fixed-width, and arbitrary-precision implementations all
recompute and certify the displayed lower bound independently, with the Lean
block computations using \texttt{native\_decide} as just stated.
\end{proof}

\begin{remark}
The same source compiled with GCC's undefined-behavior sanitizer gives the
same result.  The finer grids $1800\times400\times220$ and
$2400\times500\times280$ give the respective lower bounds
$0.033508163100$ and $0.034213110567$.  These values support convergence but
are not used in the theorem.
\end{remark}

\begin{lemma}[Explicit analytic-margin budget]\label{lem:explicit-margin}
Let $H_t$ be the envelope obtained from the displayed branch ledger by
using the shifted parameters in \eqref{eq:margin-parameters} with margin
$t$ on $7/6<\alpha<61/50$; outside that retained range set $H_t=H$.  Put
\[
 \Delta_t:=\int_{7/6}^{5/4}\max\{0,4\alpha-H_t(\alpha)\}\,d\alpha.
\]
For $0\le t\le10^{-20}$,
\begin{equation}\label{eq:explicit-margin-loss}
 |\Delta_t-\Delta|\le10^{11}t.
\end{equation}
Consequently the fixed choice
\begin{equation}\label{eq:explicit-nu}
 \nu_0=10^{-20}
\end{equation}
satisfies all interior inequalities and spends less than half of the exact
endpoint margin:
\[
 10^{11}\nu_0=10^{-9}<\frac1{10000}
 <\frac{2994511739}{18000000000000}.
\]
\end{lemma}

\begin{proof}
Only $7/6<\alpha<61/50$ contributes to the certificate.  On this interval
\[
 \gamma\ge\frac1{25},\qquad \tau\ge\frac7{50},\qquad
 \frac12-\beta_1-\beta_2\ge\frac3{50}
 \quad(\tau<\beta_1<a,\ \beta_2<\beta_1).
\]
Hence, for $0\le t\le10^{-20}$,
\[
 \gamma_t>\frac1{50},\qquad \tau_t>\frac1{10},\qquad
 \frac12-\beta_1-\beta_2-t>\frac1{20}.
\]
All integration variables are therefore at least $1/50$.  Moreover
$|\log(\tau_t/\gamma_t)|<\log4<2$, $0\le W_3\le1$, and
$0\le U_2\le U_{{\rm LS},t}\le40$.  The absolute densities in the
three-prime and two-prime integrals are consequently bounded by
\[
 \frac{61}{50}\,50^4<8\cdot10^6,\qquad
 \frac{61}{50}\,40\,50^2<2\cdot10^5,
\]
respectively; the elementary base term in \eqref{eq:H} is below $400$.
These are explicit integrable majorants on a unit box.

Away from a branch boundary, differentiate the rational branches with
respect to $t$.  Since $|\Phi'|\le1$, $|d s_{3,t}/dt|=1/\beta_3\le50$,
and
\[
 \left|\frac d{dt}\frac2{1/2-\beta_1-\beta_2-t}\right|
 \le800,
\]
write
\[
 L_t=\log(\tau_t/\gamma_t),\qquad
 Q(L)=u_+-u_-L+\frac{u_+}{2}L^2.
\]
The bounds above give $|Q(L_t)|\le5$, $|Q'(L_t)|\le3$, and
\[
 |L_t'|\le\frac{3/2}{\tau_t}+\frac4{\gamma_t}\le215.
\]
Consequently
\[
 \left|\frac d{dt}\frac{\alpha Q(L_t)}{\gamma_t}\right|
 \le\frac{4\alpha}{\gamma_t^2}|Q(L_t)|
   +\frac{\alpha}{\gamma_t}|Q'(L_t)L_t'|<10^6.
\]
On a fixed Type~II indicator class, the derivative of the three-prime
density is at most
\[
 \frac{61}{50}\,50^4\cdot50<4\cdot10^8.
\]
For the recursive candidate, Leibniz' rule and
$|\partial_tW_{3,t}|\le1/\beta_3$ give
\[
 |\partial_t C_{{\rm rec},t}|
 \le\frac4{\gamma_t^2}
   +\frac{4}{\gamma_t^2}
   +\int_{\gamma_t}^{\beta_2}\frac{d\beta_3}{\beta_3^3}
 <3\cdot10^4.
\]
Together with the $800$ bound for $U_{{\rm LS},t}$, the direct Type~II
candidate (which is independent of $t$), and
$\alpha/(\beta_1\beta_2)\le(61/50)50^2$, the complete two-prime density
varies by less than $10^8t$.  Maxima, minima, and the positive-part map are
$1$-Lipschitz when applied to this fixed family of everywhere-defined
candidates.  Lemma~\ref{lem:recursive-candidate-nonnegative} supplies the
only point that was previously conditional on candidate availability.

It remains to account for changing branch indicators and integration
limits.  There are eight hyperplanes from the four subset sums against the
two Type~II endpoints, two from the pair Type~II interval, one from the
recursive $\xi_t$ test, and five occurrences of moving $\gamma_t$ or
$\tau_t$ integration limits.  Thus even counting repeated priority
occurrences separately gives fewer than $32$ moving affine hyperplanes.
Each right-hand side moves by at most $5t$.
Because some coefficient of a moving hyperplane is $1$, its symmetric
difference inside the unit box has measure at most $10t$.  Using the
$8\cdot10^6$ majorant, all moving slabs contribute less than
\(32\cdot10t\cdot8\cdot10^6<3\cdot10^9t\).
Together with the smooth variation above and the fact that the
$\alpha$-interval has length below $1$, this is less than $10^{11}t$.
Finally, $r\mapsto\max(0,r)$ is $1$-Lipschitz, which proves
\eqref{eq:explicit-margin-loss}.

For clarity, no separate grid-limit loss is hidden here.  The canonical
mesh has $\alpha$-width $1/14400$, three-prime widths at most $1/3000$,
tail $\beta_1$-width at most $1/2000$, and containing child width at most
$9/8000$.  Proposition~\ref{prop:certificate} assigns proof-adverse values
to every one of these boxes, including boxes meeting Type~II or priority
boundaries.  Thus the locked value $32303187971$ is already the
post-discard lower bound; the only subsequent loss is the margin comparison
\eqref{eq:explicit-margin-loss}.

The affine requirements are immediate from
$4\nu_0<1/25\le\gamma$ and
$5\nu_0<1/25\le\sigma-a$.  The last displayed rational comparison is also
checked in the finite Lean audit.  Reserve the other half of the endpoint
margin and then take $x$ large enough to absorb the power-saving and $o(1)$
terms.
\end{proof}

\section{An independent lower bound for the old subtractive term}

The original decomposition retains
\begin{equation}\label{eq:F6}
 F_6=\int_{7/6}^{5/4}\alpha
 \int_{\alpha-1}^{\sigma}
 \omega\left(\frac{\alpha-\beta}{\beta}\right)
 \frac{d\beta}{\beta^2}\,d\alpha.
\end{equation}
We need a stronger bound than the conservative $0.035631$ printed in the
original paper, but it can be obtained without numerical Buchstab
integration.

\begin{lemma}\label{lem:F6}
One has
\[
 F_6>\frac{149403}{2500000}=0.0597612.
\]
\end{lemma}

\begin{proof}
Set $v=\alpha/\beta-1$.  The inner integral in \eqref{eq:F6} becomes
\[
 \int_{(4\alpha-2)/(2-\alpha)}^{1/(\alpha-1)}\omega(v)\,dv.
\]
Throughout $7/6\le\alpha\le5/4$, this interval lies in $[3.2,6]$.
Thus the lower bound in \eqref{eq:omega-bounds} applies and
\[
 F_6\ge0.5607\int_{7/6}^{5/4}w(\alpha)\,d\alpha,
 \qquad
 w(\alpha)=\frac1{\alpha-1}+4-\frac6{2-\alpha}.
\]
\par\noindent Moreover,
\[
 w''(\alpha)=\frac2{(\alpha-1)^3}-\frac{12}{(2-\alpha)^3}>0,
\]
since the first term is at least $128$ and the second at most $768/27$.
Therefore the composite midpoint sum is a lower bound.  Ten equal panels
give, by exact rational arithmetic,
\begin{align*}
 0.5607\int_{7/6}^{5/4}w(\alpha)\,d\alpha
 &\ge \frac{5607}{10000}\frac1{120}
 \sum_{j=0}^{9}w\left(\frac76+\frac{2j+1}{240}\right)\\
 &=0.0597612998251544\ldots>0.0597612.
\end{align*}
\end{proof}

\section{The endpoint budget and proof of the theorem}

Merikoski's published strict upper bounds for the other terms are
\begin{align*}
 F_1&<0.0287,& F_2&<0.08622,& F_3&<0.03107,\\
 F_4&<0.00011,& F_5&=\frac{29}{72}.
\end{align*}
The standalone $1/6$ is unchanged: it is the baseline term in Merikoski's
published treatment of the contribution with $\alpha\le7/6$, separate from
the named deficiencies $F_1$--$F_4$
\cite[PDF pp.~15--16, proof of Theorem 2]{Merikoski2020}.
Combining these with Lemma \ref{lem:F6} and the complete lower bound in Proposition
\ref{prop:certificate}, the certified zero-margin contribution through
exponent $5/4$ is strictly smaller than
\begin{align}\label{eq:core}
  C&=\frac16+\frac{287}{10000}+\frac{8622}{100000}
      +\frac{3107}{100000}+\frac{11}{100000}+\frac{29}{72}\notag\\
  &\qquad-\frac{149403}{2500000}-\frac{32303187971}{10^{12}}\notag\\
  &=\frac{5611320508261}{9000000000000}.
\end{align}

Above $5/4$, the unchanged dimension-one linear sieve contributes through a
block exponent $\varpi$
\begin{equation}\label{eq:tail}
 4\int_{5/4}^{\varpi}\alpha\,d\alpha
 =2\left(\varpi^2-\left(\frac54\right)^2\right).
\end{equation}
Choose
\[
 \varpi_*:=\frac{13231}{10000}=1.3231.
\]
Exact rational arithmetic gives
\begin{align}\label{eq:budget}
 C+2\left(\varpi_*^2-\left(\frac54\right)^2\right)
 &=\frac{8997005488261}{9000000000000}<1,\\
1-\frac{8997005488261}{9000000000000}
&=\frac{2994511739}{9000000000000}.
\end{align}
At the fixed analytic margin $\nu_0$, Lemma~\ref{lem:explicit-margin}
increases the upper bound by at most $10^{-9}$, which is less than half the
last displayed margin.  Choose $x$ large enough that all power-saving,
fundamental-lemma, and smooth-partition errors together use less than the
remaining half.  The resulting normalized total is still strictly below
$1$.

Let $P_x$ denote the greatest prime factor among the values $n^2+1$ selected
by the smooth index block.  The Chebyshev--Hooley identity used in
\cite{Merikoski2020} is
\begin{equation}\label{eq:chebyshev}
 \sum_{x<p\le P_x}|\Acal_p|\log p=X_b\log x+O(x).
\end{equation}
If $P_x\le x^{\varpi_*}$, summing the smooth modulus blocks and applying
\eqref{eq:core}, \eqref{eq:tail}, and Lemma~\ref{lem:explicit-margin} gives a
strict upper bound smaller than $X_b\log x$, contradicting
\eqref{eq:chebyshev} for all sufficiently large $x$.  Hence every
sufficiently large index block $[x,2x]$ contains an $n$ such that
\begin{equation}\label{eq:block-result}
 \Pplus(n^2+1)>x^{1.3231}.
\end{equation}

Finally put $a=1323/1000=1.323$ and
$b=13231/10000=1.3231$.  Then $b-a=1/10000$ and
$a/(b-a)=13230$.  For $x>2^{13230}$ and $n\le2x$,
\[
 n^a\le(2x)^a=2^a x^a<x^{b-a}x^a=x^b.
\]
Together with \eqref{eq:block-result}, this proves Theorem \ref{thm:main}.

\section{Formal verification and scope}

The focused Lean project contains the finite weighted Buchstab identity, its
disjoint first-bad-index iteration, the sign pattern
$A_0-A_1+A_2-A_3$, the recursive node upper/lower rules, and the
subtractive-tail inequality used in Section~\ref{sec:switching}.  Its
rational-envelope module proves the elementary lower envelope used in the
finite certificate.  The recursive-certificate module and its four block
modules independently recompute the exact outward-rounded certificate; the
value and canonical modules aggregate the block totals while keeping peak
elaboration memory bounded.  The $F_6$ module proves positivity of the
displayed second derivative, convexity, the ten-panel midpoint integral lower
bound, and the exact rational comparison with $149403/2500000$; applying this
width certificate to $F_6$ still uses the quoted Buchstab lower bound.  The
endpoint module checks the exact rational comparisons, including
$1323/1000<13231/10000$ and the strict normalized budget.  The real-power
module proves the final implication $x>2^{13230}$ and $0\le n\le2x$
imply $n^{1323/1000}<x^{13231/10000}$.  The focused project builds with
Lean 4 and Mathlib v4.22.0.

Small audit modules check affine identities for the $\nu$ margins, an
explicit rational $\nu=10^{-20}$ that satisfies those identities on
$7/6<\alpha<1.22$, a half-integer cutoff, an abstract count of
representations, one-sided box geometry, algebraic normalization of
$\alpha/\gamma$, $t+1$, $\Phi$, and $U_{\rm LS}$, the $s<1$ cutoff identity
$(2/s)/((1-\omega)\log D)=2/\log D$, the exact half-margin split
$2994511739/18000000000000$, the comparison
$10^{11}\nu<2994511739/18000000000000$, and the finite Perron budget arithmetic under
stated hypotheses
$\log((V-1/2)/U)\ge 2/(4U+1)$ and
$X^{B}\cdot 2\cdot X^{C-A'}\le X^{-A}$.
These modules construct neither the actual Perron localization nor
Rosser coefficients varying with the prefix, and they do not prove any of
the six analytic transfer gates as a whole.

The formal development treats Grimmelt--Merikoski Corollaries 7.1 and 7.2
and the standard upper and lower dimension-one linear-sieve theorem as
external analytic inputs.  Sections~\ref{sec:transfer} and
\ref{sec:switching} verify their hypotheses and prove the finite
Perron/Mellin localization on paper; these analytic arguments are not Lean
theorems.
The finite Buchstab decomposition, the integer computation, its primitive
outward-rounding rules, the elementary $F_6$ convex midpoint certificate,
the exact rational endpoint arithmetic, and the real-power dyadic conversion
are checked in Lean.  The expensive $1200\times300\times160$ block equalities
are discharged by \texttt{native\_decide}; the canonical sum, division, and
endpoint comparisons are then checked without native recomputation.  Thus
the native compiler is a trust boundary only for block values independently
cross-checked by C++ and Python; downstream arithmetic and the finite
lemmas on signs, widths, clipping, and endpoint directions are kernel checked.
None of these facts
asserts that the integer algorithm equals the analytic $H(\alpha)$ integral.  No
floating-point or proprietary computer-algebra calculation is an input to
the finite certificate.

\par\medskip
\begin{remark}[AI-assisted proof development]
The author discloses that GPT-5.6 Sol was used as a reasoning and coding
assistant in developing the core recursive Buchstab switching argument, the
proof-direction and branch ledger, the Lean certificate implementation, the
exact-integer audit scripts, and the manuscript exposition. The author
reviewed the final mathematical statements and takes responsibility for the
submitted content. GPT-5.6 Sol is an AI-assisted development tool, not an
author or an independent verifier. The Lean project checks only the finite
combinatorial, integer, rounding, elementary $F_6$, rational endpoint, and
real-power conversion components; the cited
analytic Type-I/II and linear-sieve theorems remain external inputs.  The
finite Mellin--Perron transfer is proved in the paper but not formalized in
Lean.
\end{remark}

\par\medskip
\begin{remark}[Supplementary material]
The accompanying source archive contains the Lean project and locked
toolchain, the fixed-width C++ certificate
\nolinkurl{certify_harman_recursive_tail.cpp}, and the independent
arbitrary-precision Python audit
\nolinkurl{audit_harman_recursive_tail.py}.  The command line
\texttt{1200 300 160} selects the canonical grid in both external
implementations; each prints
\texttt{saving\_lower}${}={}$\texttt{0.032303187971} and
\texttt{CERTIFIED=YES}.  No precomputed binary data are needed.  The
immutable source archive is
\url{https://doi.org/10.5281/zenodo.22140874}.
The separate script
\nolinkurl{audit_harman_recursive_branch_coverage.py} exhaustively checks
the analytic branch gates summarized in Remark~\ref{rem:branch-coverage}.
A supplied build script packages the paper, PDF, locked Lean project,
audits, and analytic ledger into one checksummed review archive.
The archive also includes a source-version ledger with theorem locations and
SHA-256 hashes for the external PDFs; the external PDFs themselves are not
redistributed.
The public verification repository is
\url{https://github.com/CGandGameEngineLearner/landau-fourth-problem-1.323}.
\end{remark}

\begin{remark}[Analytic review boundary]
The finite cross-condition localization and prefix-uniform sieve are proved
in Lemma~\ref{lem:localization} and Proposition~\ref{prop:prefix-sieve}, but
they are not formalized in Lean.  Together with the quoted Type~I/II and
linear-sieve theorems, these are the parts that require conventional
analytic-number-theory refereeing.
\end{remark}

\par\medskip
\begin{remark}[Distance from Landau's problem]
Theorem \ref{thm:main} implies infinitely often a factorization
\[
 n^2+1=pm,\qquad p>n^{1.323},\qquad m<n^{0.677+o(1)},
\]
but it neither forces $m=1$ nor overcomes the parity barrier.  It is a new
greatest-prime-factor weak form of Landau's fourth problem, not a proof of
the prime-values conjecture.
\end{remark}

\appendix
\section{Details for the prefix-uniform linear sieve}
\label{app:prefix-sieve}

This appendix records the analytic bookkeeping used in
Proposition~\ref{prop:prefix-sieve}; the main text contains the proof route
and the limitation of the argument.

\subsection{One localized prefix}

Fix one exponent cell and an ordered prefix
\(\boldsymbol q=(q_1,\ldots,q_j)\), \(j\le3\), and put
\(u=q_1\cdots q_j\).  The prefix partition is kept outside the sifted
cofactor sum; write its nonnegative weight as \(c(\boldsymbol q)\).  No
prefix-dependent cofactor cutoff is needed.  Define
\[
 a_u(m):=|\Acal_{um}|\psi_P(um)\log(um),\qquad
 |\Acal_{u,d}|:=\sum_{d\mid m}a_u(m),
\]
and, with the same weight, the model sums
\[
 |\Bcal_{u,d}|:=
 X_b\sum_{d\mid m}\frac{\rho(um)}{um}
 \psi_P(um)\log(um),\qquad X_u:=|\Bcal_{u,1}|.
\]
Write
\begin{equation}\label{eq:prefix-density}
\begin{aligned}
 \Delta(k)&:=\sum_{\ell^2+1\equiv0\pmod k}b(\ell)
              -X_b\frac{\rho(k)}k,\\
 r_u(d)&:=|\Acal_{u,d}|-|\Bcal_{u,d}|,\\
 r_P(e)&:=\sum_{e\mid k}\psi_P(k)\log k\,\Delta(k).
\end{aligned}
\end{equation}
A Rosser modulus \(d\mid P(z(\boldsymbol q))\) has all prime factors below
the least prefix prime, so \((u,d)=1\).  The changes of variables
\(k=um\) and \(e=ud\) are then exact and give the important identity
\begin{equation}\label{eq:prefix-common-discrepancy}
 r_u(d)=r_P(ud).
\end{equation}
Thus the discrepancy no longer remembers which ordered prefix represents
\(u\).

Define the coprime model sequence
\[
 |\Bcal_{u,d}^{\circ}|:=X_b
 \sum_{\substack{d\mid m\\(m,u)=1}}
 \frac{\rho(um)}{um}\psi_P(um)\log(um),qquad
 X_u^{\circ}:=|\Bcal_{u,1}^{\circ}|.
\]

\begin{lemma}[Uniform prefixed model sieve]
\label{lem:uniform-prefixed-model}
Uniformly for every ordered prefix cell with $j\le3$ and
$z=q_j$, the sequence $\Bcal_u^{\circ}$ has the same dimension-one local
density and the same fundamental-lemma constants as the unprefixed model.
For the level-$D$ Rosser weights,
\begin{align*}
 \sum_{d\mid P(z)}\lambda_{d,\boldsymbol q}^{+}
 |\Bcal_{u,d}^{\circ}|
 &\le X_u^{\circ}V(z)
 \{F(s_{\boldsymbol q})+O((\log D)^{-1/6})\}+E_u,\\
 \sum_{d\mid P(z)}\lambda_{d,\boldsymbol q}^{-}
 |\Bcal_{u,d}^{\circ}|
 &\ge X_u^{\circ}V(z)
 \{f(s_{\boldsymbol q})-O((\log D)^{-1/6})\}-E_u,
\end{align*}
where the constants are independent of the exact prefix and
\(\sum_{\boldsymbol q}c(\boldsymbol q)E_u=o(\mathfrak M(P))\).
Replacing $\Bcal_u^{\circ}$ by the full $\Bcal_u$ changes the summed model
by $O(x^{-\gamma_\nu+o(1)}\mathfrak M(P))$.
\end{lemma}

\begin{proof}
On $(m,u)=1$, multiplicativity gives
\[
 \frac{\rho(um)}{um}=\frac{\rho(u)}u\frac{\rho(m)}m.
\]
The factor $\rho(u)/u$ is outside the cofactor sieve.  The Euler factor of
the remaining weight $\rho(m)/m$ gives, for odd primes,
\[
 g(p)=
 \frac{\sum_{r\ge1}\rho(p^r)p^{-r}}
      {\sum_{r\ge0}\rho(p^r)p^{-r}}
 =\begin{cases}2/(p+1),&p\equiv1\pmod4,\\0,&p\equiv3\pmod4,
 \end{cases}
\]
with $p=2$ absorbed into the fixed initial factor.  Since every prime
dividing $u$ is at least $z=q_j$, no Euler factor with $p<z$ is removed.
Thus $g_u(p)=g(p)$ for $p<z$ and
\[
 V_u(z)=\prod_{p<z}(1-g_u(p))=\prod_{p<z}(1-g(p))=V(z).
\]
The prime number theorem in progressions modulo $4$ gives, uniformly for
$2\le w<z$,
\[
 \prod_{w\le p<z}(1-g(p))^{-1}
 \le\frac{\log z}{\log w}
 \left(1+O\!\left(\frac1{\log w}\right)\right),
\]
which is the dimension-one product hypothesis used in
\cite[Theorem 11.12 and Chapter 12]{FriedlanderIwaniec2010}.

It remains to check uniformity of the smooth family.  Put $Y=P/u$ and
\(W_u(m)=\psi_P(um)\log(um)\).  On $m\asymp Y$,
\[
 Y^r|W_u^{(r)}(m)|\ll_r\log P
\]
for every fixed $r$, with a constant independent of $u$; after extracting
the harmless factor $\log P$, these are exactly the fixed seminorms in
Merikoski's proof of Lemma~8.  Hence Theorem~11.12 and the Chapter~12
fundamental lemma give the two displayed estimates with constants uniform
in the exact value of $q_j$ inside its dyadic cell.

Finally, if $(m,u)>1$, then $q_i\mid m$ for one of at most three primes
$q_i\ge x^{\gamma_\nu}$.  Writing $m=q_i m'$ and using
$\rho(n)\ll\tau(n)$ gives, after summing the ordered prefix cells,
\[
 \sum_{\boldsymbol q}c(\boldsymbol q)
 \sum_{i\le j}\sum_{q_i\mid m}
 X_b\frac{\rho(um)}{um}\psi_P(um)\log(um)
 \ll x^{-\gamma_\nu+o(1)}\mathfrak M(P).
\]
This proves the final assertion and the lemma.
\end{proof}

Absorb the negligible coprimality correction from
Lemma~\ref{lem:uniform-prefixed-model} into $E_u$ and return to the notation
$X_u=|\Bcal_{u,1}|$.
Let \(\lambda_{d,\boldsymbol q}^{\pm}\) be the level-\(D\) Rosser weights.
They are supported on squarefree \(d\mid P(z(\boldsymbol q))\), \(d<D\),
and satisfy \(|\lambda_{d,\boldsymbol q}^{\pm}|\le1\).
The Rosser inequalities and the preceding model evaluation give, with
\(E_u\ge0\) and
\(\sum_{\boldsymbol q}c(\boldsymbol q)E_u=o(\mathfrak M(P))\),
\cite[Chapter 12, (12.12)--(12.14), pp. 235--236]{FriedlanderIwaniec2010}
\begin{align}
 S(\Acal(P)_u,z(\boldsymbol q))
 &\le X_uV_u(z)
 \{F(s_{\boldsymbol q})+O((\log D)^{-1/6})\}
 +R_{\boldsymbol q}^{+}+E_u,\label{eq:prefix-upper-sieve}\\
 S(\Acal(P)_u,z(\boldsymbol q))
 &\ge X_uV_u(z)
 \{f(s_{\boldsymbol q})-O((\log D)^{-1/6})\}
 +R_{\boldsymbol q}^{-}-E_u.\label{eq:prefix-lower-sieve}
\end{align}
where \(s_{\boldsymbol q}=\log D/\log z(\boldsymbol q)\) and
\[
 R_{\boldsymbol q}^{\pm}
 =\sum_{\substack{d\mid P(z(\boldsymbol q))\\d<D}}
 \lambda_{d,\boldsymbol q}^{\pm}r_P(ud).
\]
The signs of these remainders are retained: no absolute value is taken
before the prefix sum is collected.  The upper formula is used directly for
\(s_{\boldsymbol q}\ge1\); the lower formula is used for
\(s_{\boldsymbol q}\ge2\), and zero is used below \(2\).

\subsection{Collecting the prefix sum}

For a cell with \(u\le x^{s_0}\), take
\[
 D=x^{1/2-s_0-\nu},\qquad ud\le x^{1/2-\nu}.
\]
Let \(c(\boldsymbol q)\ge0\) denote the smooth dyadic/exponent partition of
the ordered prefix; its size and all partial-summation losses are
\(O((\log x)^B)\).  Multiply \eqref{eq:prefix-upper-sieve} or
\eqref{eq:prefix-lower-sieve} by \(c(\boldsymbol q)\), sum first, and only
then take an absolute value.  Using
\eqref{eq:prefix-common-discrepancy} and collecting \(e=ud\) gives the exact
identity
\begin{equation}\label{eq:collected-prefix-remainder}
 \mathcal R^{\pm}
 :=\sum_{\boldsymbol q}c(\boldsymbol q)R_{\boldsymbol q}^{\pm}
 =\sum_{e\le x^{1/2-\nu}}\Lambda_e^{\pm}r_P(e),\qquad
 \Lambda_e^{\pm}:=
 \sum_{\substack{\boldsymbol q,d:\,ud=e\\
                   d\mid P(z(\boldsymbol q)),\ d<D}}
 c(\boldsymbol q)\lambda_{d,\boldsymbol q}^{\pm}.
\end{equation}
If an auxiliary smooth product cutoff is introduced, Mellin inversion first
writes it as a polylogarithmic-\(L^1\) superposition of the same identity
with a common smooth function of \(k/P\); Corollary~7.1 is uniform for each
such function.  Hence it creates only the displayed logarithmic loss, not a
prefix-dependent discrepancy inside the \(k\)-sum.

For fixed \(e\), the at most three ordered prefix primes divide \(e\), and
\(d=e/u\) is determined.  Hence fixed divisor bounds for \(c(\boldsymbol q)\)
and \(|\lambda_{d,\boldsymbol q}^{\pm}|\le1\) imply
\begin{equation}\label{eq:Lambda-divisor-bound}
 |\Lambda_e^{\pm}|\ll(\log x)^B\tau_K(e)
\end{equation}
with fixed \(B,K\).  Put
\[
 \delta_0=(1-2\theta)\nu-\varepsilon_{\rm GM}/4>0.
\]
Split \([P,4P]\) into a fixed number of smooth dyadic blocks and remove
\(\log k\) by partial summation.  Corollary~7.1 of
\cite{GrimmeltMerikoski2025}, with \(a=h=1\), \(\eta=\nu\), and
\(P\le x^2\), has inside its modulus-$e$ absolute value exactly
\[
 \sum_{k\equiv0\pmod e}\psi_1(k/K)
 \left\{\sum_{\ell^2+1\equiv0\pmod k}\psi_2(\ell/x)
       -\frac{\rho(k)}k\,x\widehat\psi_2(0)\right\}.
\]
Rescaling $b$ supplies $\psi_2$, the fixed decomposition of $\psi_P$
supplies $\psi_1$, and partial summation supplies $\log k$; hence this is
precisely $r_P(e)$ in \eqref{eq:prefix-density}.  Corollary~7.1 therefore
gives
\[
 \sum_{e\le x^{1/2-\nu}}|r_P(e)|
 \ll x^{1-\delta_0}(\log x)^{B_0}.
\]
Choosing \(\kappa<2\delta_0\) in
\(\tau_K(e)\ll e^\kappa\), and absorbing the dyadic logarithms, yields
\begin{equation}\label{eq:prefix-remainder-power-saving}
\begin{aligned}
 |\mathcal R^{\pm}|
 &\le\sum_{e\le x^{1/2-\nu}}|\Lambda_e^\pm|\,|r_P(e)|\\
 &\ll(\log x)^Bx^{\kappa/2}
       \sum_{e\le x^{1/2-\nu}}|r_P(e)|\\
 &\ll(\log x)^{B+B_0}x^{1-\delta_0+\kappa/2}
 \ll x^{1-\delta}.
\end{aligned}
\end{equation}
for some \(\delta=\delta(\nu)>0\).
Thus only the total multiple sum, not each prefix separately, has the
required power-saving error.

\subsection{The two sieve branches}

For the \(A_3\) lower branch, \(u=q_1q_2q_3\), \(z=q_3\), and
\[
\begin{aligned}
 s_{3,\nu}
 &=\frac{1/2-\beta_1-\beta_2-\beta_3-\nu}{\beta_3},\\
 \widehat s_{3,\nu}
 &=\frac{1/2-s_0-\nu}{\beta_3}+O(1/\log x)
 \le s_{3,\nu}+O(1/\log x).
\end{aligned}
\]
The cell upper endpoint therefore gives a lower sieve weight.  Replace
\(e^{-\gamma_E}f(\widehat s_{3,\nu})\) by
\(\Phi(\widehat s_{3,\nu})\); if \(\widehat s_{3,\nu}\le2\), use zero.
At the endpoint \(\alpha=7/6\),
\[
 3\tau_\nu=\frac12-\frac92\nu,\qquad
 \frac12-s_0-\nu>\frac72\nu,
\]
so \(D\) remains a positive power of \(x\).  This proves the three-prime
remainder is admissible at the fixed margin \(\nu=\nu_0\).

For the two-prime upper branch, \(u=q_1q_2\), \(z=q_2\), and
\[
 \widehat s_{2,\nu}
 =\frac{1/2-s_0-\nu}{\beta_2}+O(1/\log x)>0.
\]
If \(1\le\widehat s_{2,\nu}\le3\), then
\(e^{-\gamma_E}F(\widehat s_{2,\nu})=2/\widehat s_{2,\nu}\);
if \(\widehat s_{2,\nu}\ge3\), use \(2/(3\beta_2)\).
When \(0<\widehat s_{2,\nu}<1\), put
\(z_*=D^{1-\omega}<D<q_2\).  Positivity gives
\[
 S(\Acal(P)_{q_1q_2},q_2)
 \le S(\Acal(P)_{q_1q_2},z_*),
\]
whose sieve parameter is \(1/(1-\omega)>1\).  For fixed \(\omega>0\)
the upper-sieve formula applies, and
\[
 \frac{e^{-\gamma_E}F(1/(1-\omega))}{\log z_*}
 =\frac{2}{\log D}.
\]
After \(x\to\infty\), letting \(\omega\to0^+\) therefore gives the first
term in
\[
 U_{{\rm LS},\nu}(\beta_1,\beta_2)
 =\max\left\{\frac2{1/2-\beta_1-\beta_2-\nu},
              \frac2{3\beta_2}\right\}.
\]
The finite mesh uses the proof-adverse endpoints; boxes without positive
level contribute no certified tail saving.

Finally, for \(z(\boldsymbol q)\in[Z,2Z]\), \(Z=x^\zeta\),
\[
 \frac{\log D}{\log z(\boldsymbol q)}
 =\frac{1/2-s_0-\nu}{\zeta}+O(1/\log x).
\]
The functions \(F,f\) are continuous at the transition points, and the
certificate uses the lower value for the lower branch and the upper value
for the upper branch.  Hence freezing the multiplier on a cell changes the
total model sum by \(o(\mathfrak M(P))\).  Merikoski's Lemma~7 and the prime
number theorem then give, on an \(A_3\) cell,
\[
 (1+o(1))\mathfrak M(P)\,\alpha
 \int_{\rm cell}e^{-\gamma_E}f(s_{3,\nu})
 \frac{d\beta_1d\beta_2d\beta_3}
      {\beta_1\beta_2\beta_3^2}.
\]
Since \(\Phi(s)\le e^{-\gamma_E}f(s)\), summing the cells gives
\eqref{eq:A3lower}.  For the two-prime upper branch the corresponding
coefficient relative to
\(d\beta_1d\beta_2/(\beta_1\beta_2)\) is
\(e^{-\gamma_E}F(s_{2,\nu})/\beta_2\).  It equals
\(2/(1/2-\beta_1-\beta_2-\nu)\) for \(1\le s_{2,\nu}\le3\), is at most
\(2/(3\beta_2)\) for \(s_{2,\nu}\ge3\), and has the first bound after the
cutoff reduction when \(0<s_{2,\nu}<1\).  This is precisely
\(U_{{\rm LS},\nu}\), and summing gives its term in \eqref{eq:Tlower}.
All fundamental-lemma and Type-I estimates used here are the cited external
analytic inputs; the prefix collection itself is the exact calculation
\eqref{eq:collected-prefix-remainder}.


\begin{thebibliography}{99}

\bibitem{Carella2025}
N. A. Carella,
\emph{Largest prime factors of polynomials},
arXiv:2308.10075v7, 2025.

\bibitem{GrimmeltMerikoski2025}
L. Grimmelt and J. Merikoski,
\emph{On the greatest prime factor and uniform equidistribution of quadratic polynomials},
arXiv:2505.00493v2, 2025.

\bibitem{FriedlanderIwaniec2010}
J. Friedlander and H. Iwaniec,
\emph{Opera de Cribro},
American Mathematical Society Colloquium Publications, vol.~57,
American Mathematical Society, Providence, 2010.

\bibitem{Iwaniec1978}
H. Iwaniec,
\emph{Almost-primes represented by quadratic polynomials},
Invent. Math. \textbf{47} (1978), 171--188,
\href{https://doi.org/10.1007/BF01578070}{doi:10.1007/BF01578070}.

\bibitem{Li2025}
R. Li,
\emph{On the largest prime factor of quadratic polynomials},
arXiv:2406.07575v2, 2025.

\bibitem{Merikoski2020}
J. Merikoski,
\emph{On the largest prime factor of $n^2+1$},
J. Eur. Math. Soc. (JEMS) \textbf{25} (2023), 1253--1284,
\href{https://doi.org/10.4171/JEMS/1216}{doi:10.4171/JEMS/1216}.

\bibitem{Pascadi2026}
A. Pascadi,
\emph{Large sieve inequalities for exceptional Maass forms and the greatest prime factor of $n^2+1$},
Forum Math. Pi \textbf{14} (2026), e8,
1--54,
\href{https://doi.org/10.1017/fmp.2026.10025}{doi:10.1017/fmp.2026.10025}.

\end{thebibliography}

\end{document}
